\section{Cái trần cao bao nhiêu}
\label{sec:ch13-tran-cao-bao-nhieu}

\index{giới hạn lý thuyết!độ cao theo lớp hàm|(}%
Định lý~\ref{thm:ch13-khoang-cach-do-phuc-tap} nói rằng ô ở góc hình là rỗng
nhưng không nói nó rỗng tới đâu, vì nó không giả thiết gì về $f$. Muốn biết cái
trần cao bao nhiêu thì phải nói $f$ thuộc lớp hàm nào, và mục này đọc các chặn
mà bài rút ra cho từng lớp. Trật tự đi từ trường hợp tệ nhất tới trường hợp có
lối ra.

\index{hàm ngẫu nhiên!không giải thích được}%
Trường hợp tệ nhất là một hàm không có cấu trúc nào cả. Lấy $\mathcal{X}$ là tập
mọi chuỗi $n$ bit, $\mathcal{Y} = \{0,1\}$, và lấy $f$ ngẫu nhiên đều trong số
mọi hàm Boole trên $\mathcal{X}$. Bài chứng minh rằng với một hằng số
$0 < \epsilon < 1$ cố định và mọi ngân sách $b \le (1-\epsilon)2^n$, hầu như
chắc chắn theo phép chọn $f$, mọi hàm $g$ với $K(g) \le b$ đều sai trên gần đúng
một nửa số đầu vào~\autocite{p26limits}. Nói cách khác, một lời giải thích không
dùng gần trọn $2^n$ bit thì không hơn gì tung đồng xu.

Chứng minh là một phép đếm và đáng viết ra vì nó cho thấy cái trần này đến từ
đâu. Với $g$ cố định, mỗi đầu vào $x$ cho $f(x) = g(x)$ với xác suất một nửa và
các đầu vào độc lập nhau, nên số chỗ trùng nhau tuân theo phân phối nhị thức và
vượt quá một nửa một lượng cố định nào đó với xác suất nhỏ theo hàm mũ của
$2^n$, theo chặn Chernoff. Số hàm $g$ phải xét thì không quá $2^{b+1}$, đúng
phép đếm của
mục~\ref{sec:ch13-ham-xap-xi}. Nhân hai đại lượng ấy lại theo chặn hợp: khi $b$
nhỏ hơn $2^n$ một phần hằng số, số hàm ứng viên tăng chậm hơn xác suất từng hàm
trượt giảm, nên tổng vẫn nhỏ. Cái trần ở đây không phải một tính chất tinh vi
của hàm ngẫu nhiên, mà là chuyện có quá ít chương trình ngắn so với số hàm cần
mô tả.

Phải đặt kết quả ấy đúng chỗ, vì nó dễ bị đọc quá tay. Một mô hình đã huấn luyện
không phải một hàm Boole ngẫu nhiên: nó được chọn ra bằng một quy trình tối ưu
trên dữ liệu có cấu trúc, và bản thân trọng số của nó là một mô tả hữu hạn, nên
$K(f)$ của nó nhỏ hơn hẳn $2^n$. Bài không tuyên bố ngược lại; câu bài viết là
kết quả này có ý nghĩa với những hệ thống AI phức tạp mà hành vi có thể giống
hàm ngẫu nhiên ở một số vùng của không gian đầu vào~\autocite{p26limits}. Phát
biểu ấy là một phỏng đoán về mô hình thật chứ không phải một hệ quả của định lý,
và bài không đưa bằng chứng nào cho nó.

\index{hằng số Lipschitz}%
Lớp hàm thứ hai gần thực tế hơn nhiều. Một hàm $f$ trên $[0,1]^d$ gọi là
$L$-Lipschitz nếu $|f(x) - f(y)| \le L\|x-y\|_2$ với mọi $x, y$, tức tốc độ đổi
đầu ra bị chặn bởi hằng số $L$. Nhiều mô hình học máy thỏa tính chất này, trong
đó có mạng nơ-ron với trọng số bị chặn và hàm kích hoạt thích
hợp~\autocite{p26limits}. Với lớp này bài dựng được chặn theo cả hai chiều.

Chiều trên đến từ một phép dựng thẳng. Chia $[0,1]^d$ thành lưới đều gồm $m^d$
hình hộp cạnh $1/m$, rồi lấy $g$ nhận giá trị $f$ tại tâm hộp trên toàn hộp ấy.
Khoảng cách từ một điểm bất kỳ tới tâm hộp chứa nó không quá $\sqrt{d}/(2m)$,
nên theo tính Lipschitz, sai số không quá $L\sqrt{d}/(2m)$; muốn nó dưới $\delta$
thì lấy $m$ cỡ $L/\delta$. Mô tả $g$ tốn hai phần: độ phân giải lưới, tốn
$\mathcal{O}(\log(L/\delta))$ bit, và một giá trị cho mỗi tâm hộp, tức $m^d$ giá
trị, mỗi giá trị cần $\mathcal{O}(\log(L/\delta))$ bit độ chính xác. Cộng lại:
\[
  \kappa_f(\delta) \;=\; \mathcal{O}\!\left( \left(\frac{L}{\delta}\right)^{\!d} \log \frac{L}{\delta} \right).
\]

\index{lời nguyền số chiều}%
Đọc công thức ấy theo hai biến thì ra hai câu khác nhau hẳn. Theo $\delta$, độ
phức tạp tăng theo đa thức: muốn sai số nhỏ đi mười lần thì trả một cái giá chịu
được. Theo $d$, nó tăng theo hàm mũ, và bài gọi đó là \tn{lời nguyền số
chiều}{curse of dimensionality} đặt vào bài toán giải thích, cùng cái tên mà
mục~\ref{sec:ch07-dieu-chua-giai-quyet} đã gặp ở chỗ khác~\autocite{p26limits}.
Với một mô hình nhận đầu vào nghìn chiều, số hạng $(L/\delta)^d$ lớn tới mức
không có nghĩa vật lý nào.

Chiều dưới hỏi ngược lại: có phải mọi lời giải thích đều phải trả cái giá ấy,
hay chỉ phép dựng lưới mới phải trả.

\begin{theorem}[Chặn dưới cho hàm Lipschitz]
\label{thm:ch13-chan-duoi-lipschitz}
Với mọi số chiều $d \ge 1$, mọi hằng số $L > 0$ và mọi ngân sách
$b \in \mathbb{N}$, tồn tại một hàm $L$-Lipschitz $f : [0,1]^d \to \mathbb{R}$
sao cho mọi lời giải thích $g$ với $K(g) \le b$ đều có
$\mathcal{E}_{\infty}(f,g) = \Omega\big(L \cdot 2^{-b/d}\big)$.
\end{theorem}

Định lý~\ref{thm:ch13-chan-duoi-lipschitz} là định lý 3.4 của
bài~\autocite{p26limits}. Sách giữ phát biểu ấy nhưng không giữ chứng minh của
bài, nên mục này đọc chứng minh ấy trước, rồi dựng một chứng minh khác.

Ý của bài đi theo hình học. Một hàm $g$ trong ngân sách $b$ được cho là chia
$[0,1]^d$ thành không quá $2^{b+1}$ vùng nhận giá trị khác nhau, nên theo nguyên
lý ngăn kéo có ít nhất một vùng thể tích không dưới $2^{-(b+1)}$; một vùng thể
tích ấy trong $\mathbb{R}^d$ có đường kính không dưới cỡ $2^{-b/d}$ theo bất
đẳng thức đẳng chu; chọn $f$ biến thiên với độ dốc đúng bằng $L$ dọc theo đường
kính ấy thì $f$ đổi một lượng cỡ $L \cdot 2^{-b/d}$ trong khi $g$ giữ một giá
trị hằng.

Hai bước trong lập luận ấy không đứng được.

Bước thứ nhất là phép đếm. Con số $2^{b+1}$ đếm \emph{số hàm} có độ phức tạp
không quá $b$, đúng như bổ đề mà mục~\ref{sec:ch13-ham-xap-xi} đã dùng. Nó không
nói gì về số vùng giá trị của một hàm trong số ấy. Hàm đồng nhất $g(x) = x$ có
mô tả dài đúng một hằng số và nhận vô hạn giá trị khác nhau trên $[0,1]$: ngân
sách nhỏ không buộc một hàm phải là hàm hằng từng khúc. Chính con số $2^{b+1}$
đã xuất hiện ở đầu mục này và ở đó nó đúng, vì nó chặn số ứng viên $g$ trong một
chặn hợp. Chỗ hụt là mang con số ấy sang làm số mảnh của một phép chia miền.

Bước thứ hai là thứ tự lượng từ. Phát biểu nói \emph{tồn tại} một $f$ sao cho
\emph{mọi} $g$ trong ngân sách đều sai nhiều. Lập luận lại cố định $g$ trước, rồi
dựng $f$ từ vùng mà chính $g$ ấy sinh ra và từ đường kính của vùng đó. Cái nhận
được là \enquote{với mọi $g$ tồn tại một $f$}, yếu hơn và không cho định lý.

Phát biểu thì vẫn đúng, và một phép đếm khác cho nó. Chia $[0,1]^d$ thành $m^d$
ô cạnh $1/m$, và trong mỗi ô $c$ đặt một cái chóp
\[
  \varphi_c(x) \;=\; \max\!\big(0,\; h - L\lVert x - t_c \rVert_2\big),
  \qquad h = \frac{L}{2m},
\]
với $t_c$ là tâm ô. Mỗi chóp là hàm $L$-Lipschitz, và giá của nó là quả cầu bán
kính $h/L = 1/(2m)$ quanh tâm, nằm gọn trong ô, nên phần trong của các giá rời
nhau. Với mỗi
vector dấu $s \in \{-1,+1\}^{m^d}$ đặt $f_s = \sum_c s_c \varphi_c$.

Mỗi $f_s$ vẫn là $L$-Lipschitz, và chỗ này đáng viết ra vì nó không phải hệ quả
của riêng việc các giá rời nhau. Lấy hai điểm $u$ và $v$. Nếu cả hai nằm trong
cùng một giá thì chỉ một cái chóp khác không ở đó và bất đẳng thức là bất đẳng
thức của chóp ấy. Nếu chúng nằm ở hai giá khác nhau thì đoạn thẳng nối chúng cắt
qua vùng mà mọi chóp đều bằng không, nên
$|f_s(u) - f_s(v)| \le |f_s(u)| + |f_s(v)|$, và mỗi số hạng lại không vượt quá
$L$ nhân khoảng cách từ điểm ấy tới biên giá của nó. Cộng hai khoảng cách ấy lại
thì không vượt quá $\lVert u - v \rVert_2$. Điều làm cho lập luận chạy là chiều
cao $h$ đã chọn: nó đúng bằng $L$ nhân bán kính giá, nên chóp chạm không đúng
lúc chạm biên ô.

Hai hàm khác dấu ở ô $c$ thì lệch nhau $2h = L/m$ ngay tại
tâm ô ấy, nên
\[
  \lVert f_s - f_{s'} \rVert_{\infty} \;\ge\; \frac{L}{m}
  \qquad \text{khi } s \ne s'.
\]

Họ ấy có $2^{m^d}$ hàm. Giả sử mọi hàm $L$-Lipschitz đều được một $g$ nào đó với
$K(g) \le b$ xấp xỉ trong sai số $\delta$. Nếu một $g$ phục vụ được cả $f_s$ lẫn
$f_{s'}$ thì $\lVert f_s - f_{s'} \rVert_{\infty} \le 2\delta$, nên khi
$2\delta < L/m$ hai vector dấu khác nhau buộc phải dùng hai $g$ khác nhau. Lấy
$\delta \le L/(3m)$ thì phép gán $f_s \mapsto g$ là đơn ánh, do đó
$2^{m^d} \le 2^{b+1}$, tức $m^d \le b+1$. Chọn
$m = \lfloor L/(3\delta) \rfloor$ và đọc lại theo $\delta$:
\[
  \kappa_f(\delta) \;\ge\; \Big\lfloor \frac{L}{3\delta} \Big\rfloor^{d} - 1
  \;=\; \Omega\!\left( \left(\frac{L}{\delta}\right)^{\!d} \right)
\]
cho một hàm $f$ nào đó trong họ vừa dựng. Thứ tự lượng từ lần này đúng ngay
trong cách lập luận: phép đếm chỉ thẳng ra rằng \emph{có} một $f$ đánh bại
\emph{mọi} $g$ trong ngân sách.

Chặn vừa dựng mạnh hơn chặn của bài, vì $b^{-1/d}$ lớn hơn hẳn $2^{-b/d}$, và nó
kéo theo phát biểu của định lý~\ref{thm:ch13-chan-duoi-lipschitz}. Đó là lý do
sách giữ định lý và bỏ chứng minh.

Hệ quả cho chương là cụ thể, và phải phát biểu đúng dạng lượng từ. Ở ngân sách
$b$, phép dựng lưới đạt sai số cỡ
$L \cdot b^{-1/d}$ nhân thêm một thừa số lôgarit, vì $b$ bit mua được cỡ $b$ ô
lưới sau khi đã trả tiền mã hóa giá trị tại mỗi ô, và điều đó đúng với
\emph{mọi} hàm $L$-Lipschitz. Phép đếm vừa rồi đi chiều ngược lại và chỉ nói về
\emph{một} hàm: trong họ hàm chóp có một hàm mà không lời giải thích nào trong
ngân sách hạ được sai số xuống dưới cỡ $L \cdot b^{-1/d}$. Ghép hai chiều lại
thì trường hợp xấu nhất của lớp Lipschitz có sai số đúng cỡ $L \cdot b^{-1/d}$,
sai khác một thừa số lôgarit và một hằng số phụ thuộc $d$.

Đọc theo $\delta$ thì câu gọn hơn: có những hàm $L$-Lipschitz mà giải thích
chúng trong sai số $\delta$ tốn ít nhất cỡ $(L/\delta)^d$ bit. Vậy lời
nguyền số chiều ở đây không phải một tính chất riêng của phép dựng lưới, mà là
tính chất của chính bài toán trong trường hợp xấu nhất. Nó không nói mọi hàm
Lipschitz đều đắt như thế, và đoạn cuối mục này quay lại đúng chỗ ấy.

Nhận xét 3.5 của bài kết luận đúng chỗ này: hai chặn khớp nhau, và sự phụ thuộc
hàm mũ vào số chiều là không tránh được trong trường hợp tổng quát~\autocite{p26limits}.
Cái bài không có là con đường tới kết luận ấy. Nhận xét đọc nó ra từ định lý 3.4
của bài, tức định lý~\ref{thm:ch13-chan-duoi-lipschitz} ở trên, vốn quá yếu để
đỡ được, và mệnh đề 3.6 thì viết chặn trên thành
$\mathcal{O}(L \cdot 2^{-b/d})$ bằng cách cho rằng ngân sách $b$ mua được $2^{b}$
ô, tức bỏ qua chính khoản tiền mã hóa giá trị mà chứng minh chặn trên đã tính.
Bước ấy sai, và sách không dùng nó; điều sách nhận là kết luận, với phép đếm ở
trên làm chỗ dựa. Hệ quả 3.8 ở cuối mục này hụt theo cùng một kiểu và không được
cứu như vậy.

Chú ý cả hai chặn dưới đều có dạng tồn tại. Chúng nói có những hàm Lipschitz
khó tới mức đó, không nói mọi hàm Lipschitz đều thế. Hàm hằng và hàm tuyến tính
là hàm Lipschitz và giải thích được rất rẻ; bài ghi rõ điều này ngay sau định
lý~\autocite{p26limits}.

\index{độ phức tạp Kolmogorov!xấp xỉ cho từng lớp mô hình}%
Còn một khoảng cách nữa phải bắc cầu trước khi các chặn này chạm được thực
hành. \tn{Độ phức tạp Kolmogorov}{Kolmogorov complexity} không tính
được, nên $\kappa_f$ và
$\varepsilon_f$ không tính được. Bài xử lý bằng cách thay $K$ bằng độ dài mã hóa
của từng lớp mô hình cụ thể: với mỗi lớp, đếm thông tin tối thiểu cần để chỉ
định một phần tử của lớp ấy. Bảng~\ref{tab:ch13-do-phuc-tap-lop-mo-hinh} chép
năm kết quả của phép đếm đó.

\begin{table}[htbp]
  \centering
  \caption{Năm ước lượng mà bài dùng thay cho $K(g)$, chép từ mệnh đề 3.7 của
  bài báo~\autocite{p26limits}. Bài gọi lập luận của mình là phác thảo chứng
  minh chứ không phải chứng minh. Ở dòng cuối, bài đếm riêng trọng số và độ
  chệch, viết chữ cái thứ hai là $b$; sách gộp chúng vào $w$ vì $b$ ở chương này
  đã là ngân sách bit.}
  \label{tab:ch13-do-phuc-tap-lop-mo-hinh}
  \begin{tabular}{@{}llp{5cm}@{}}
    \toprule
    Lớp mô hình & Ước lượng $K(g)$ & Phụ thuộc vào \\
    \midrule
    Mô hình tuyến tính & $\mathcal{O}(n \log n)$ & $n$ \tn{đặc trưng}{feature} \\
    Cây quyết định & $\mathcal{O}(|T| \log |T|)$ & $|T|$ nút \\
    Danh sách luật & $\mathcal{O}(r\,l \log (r\,l))$ & $r$ luật, độ dài trung bình $l$ \\
    Láng giềng gần nhất & $\mathcal{O}(m\,d \log (m\,d))$ & $m$ mẫu lưu, $d$ chiều \\
    Mạng nơ-ron & $\mathcal{O}(w \log p + a)$ & $w$ tham số, mô tả kiến trúc dài $a$, độ chính xác $p$ bit \\
    \bottomrule
  \end{tabular}
\end{table}

Bảng ấy là chỗ toàn bộ phần lý thuyết đổi loại phát biểu, nên nó đáng một câu
cảnh giác. Mỗi dòng là độ dài của một cách mã hóa mà bài tự chọn, tức một chặn
trên cho $K$ của lớp ấy, không phải giá trị của $K$. Một mô hình tuyến tính mà
mọi hệ số bằng nhau có $K$ nhỏ hơn $\mathcal{O}(n \log n)$ rất nhiều. Cho nên
các chặn ở những mục sau, khi thay $K$ bằng cột giữa của bảng, chỉ giữ được
chiều nào mà một chặn trên còn dùng được, và
mục~\ref{sec:ch13-dieu-chua-giai-quyet} quay lại chỗ này.

Bài rút hai hệ quả trực tiếp từ bảng. Với cây quyết định $|T|$ nút xấp xỉ một
hàm có \tn{biến phân bị chặn}{bounded variation} $V(f)$, sai số là
$\mathcal{O}\big(V(f) \cdot |T|^{-1/d}\big)$: cây càng lớn thì sai số càng nhỏ,
nhưng tốc độ giảm chậm đi theo hàm mũ của số chiều~\autocite{p26limits}. Với
mạng nơ-ron xấp xỉ một hàm $L$-Lipschitz tới sai số $\delta$, chặn là
$\mathcal{O}\big((L/\delta)^d \log(L/\delta) + d \log d\big)$, tức cùng một dáng
với phép dựng lưới ở trên~\autocite{p26limits}.

Bài dành một mục riêng, mục 3.3, cho cùng cái trần nhìn từ phía lý thuyết thông
tin của Shannon thay vì từ phía AIT: nó chặn sai số kỳ vọng theo khoảng cách
giữa lượng thông tin tương hỗ mà $f$ mang và lượng mà $g$ mang được, rồi nối
$\kappa_f$ với hàm tốc độ và độ méo của lý thuyết nén có mất
mát~\autocite{p26limits}. Sách không đọc mục ấy, vì hai kết quả trong đó không
đổi kết luận nào của chương và mục~\ref{sec:ch13-tran-noi-gi} lập luận trên cái
trần chứ không trên đường tới nó; câu hỏi 7 ở cuối chương chỉ đường tới đó cho
người đọc muốn đi tiếp.

Bài còn một hệ quả thứ ba, cho mô hình tuyến tính, và sách không dùng nó. Hệ quả
3.8 phát biểu rằng sai số của một mô hình tuyến tính $n$ đặc trưng xấp xỉ một
hàm $L$-Lipschitz là $\Omega\big(L \cdot d^{1/2} \cdot n^{-1/d}\big)$. Chứng
minh của nó đọc chặn dưới ấy ra từ chặn trên của phép dựng lưới, tức xử lý một
$\mathcal{O}$ như thể nó là một $\Theta$: muốn một sàn cho sai số thì phải có
một sàn cho $\kappa_f$, mà bài chỉ có một trần. Kết luận có thể vẫn đúng, nhưng
không đúng theo đường ấy, và nó cũng không đúng dưới dạng ấy, vì sai số của mô
hình tuyến tính khớp nhất với một $f$ cố định do khoảng cách từ $f$ tới lớp
tuyến tính quyết định chứ không giảm dần khi thêm đặc trưng. Chỗ nào chương này
cần một phát biểu về mô hình tuyến tính thì lấy từ định
lý~\ref{thm:ch13-chan-duoi-lipschitz}, với phép đếm ở đầu mục làm chỗ dựa, cộng
với dòng đầu của bảng~\ref{tab:ch13-do-phuc-tap-lop-mo-hinh}: gán cho một mô
hình tuyến tính $n$ đặc trưng ngân sách $b$ mà dòng ấy cho, thì tồn tại những
hàm $L$-Lipschitz mà mọi mô hình tuyến tính trong ngân sách ấy đều mắc sai số ít
nhất cỡ $L \cdot b^{-1/d}$.

\index{giả thuyết đa tạp}%
Cuối mục là lối ra, và nó là kết quả lạc quan duy nhất trong bài. Các chặn ở
trên đều tính theo $d$, số chiều của không gian đầu vào. Nhưng dữ liệu thật hiếm
khi trải đều khắp không gian ấy; nó thường tụ quanh một cấu trúc chiều thấp hơn,
điều mà giới học máy gọi là \tn{giả thuyết đa tạp}{manifold hypothesis}. Bài đo
chiều của cấu trúc ấy bằng \tn{số chiều đếm hộp}{box-counting dimension}: nếu
phủ tập giá của phân phối cần khoảng $\eta^{-d_B}$ hộp cạnh $\eta$ khi $\eta$
tiến về không, thì tập ấy có số chiều đếm hộp $d_B$. Bài chứng minh rằng khi
$d_B < d$, độ phức tạp cần cho sai số kỳ vọng $\delta$ dưới phân phối ấy là
$\mathcal{O}\big((1/\delta)^{d_B} \log(1/\delta)\big)$~\autocite{p26limits}. Số
mũ đổi từ $d$ sang $d_B$: cái quyết định giá của một lời giải thích là chiều nội
tại của dữ liệu, không phải chiều của không gian chứa nó. Bài đưa một con số
minh họa cho ý ấy: dữ liệu nghìn chiều nằm gần một đa tạp mười chiều thì số mũ
là mười chứ không phải một nghìn~\autocite{p26limits}.

\index{giới hạn lý thuyết!độ cao theo lớp hàm|)}%
